% GENERATED FILE. Edit canonical lessons, then run npm run generate:books.
% canonical-source-hash: fnv1a64:82e379162b7394d7
% canonical-lessons: JA-C405-teien, JA-C405-yuumei, JA-C405-migoto, JA-C405-areru, JA-C405-suzumu

\chapter{A Famous Garden}
\label{ch:ja-a-famous-garden}
\hlchaptermodality{405}

\begin{chapteropening}
\textbf{By the end of this chapter:} \emph{I can describe a famous garden.}

Complete the last lesson of chapter 405: I can describe a famous garden.
\end{chapteropening}

\section[\texorpdfstring{\ja{ていえん} (teien)}{ていえん (teien)}]{\ja{ていえん} (teien) — a landscaped garden}

\label{lesson:JA-C405-teien}



\begin{quote}
Before the new one: say the Japanese for wonderful, then the Japanese for to look down on.
\end{quote}

\subsection*{You'll want to know: \ja{ていえん}}
\textbf{\ja{ていえん}} — \emph{teien} — \textquotedblleft{}a landscaped garden\textquotedblright{}.

\begin{quote}
\textbf{\ja{てらの} \ja{ていえんは} \ja{しずかでした。}} — \emph{tera no teien wa shizuka deshita} — \textquotedblleft{}The temple garden was quiet.\textquotedblright{}
\end{quote}

\begin{cousinweb}
\textbf{\ja{ていえん}} (\emph{teien}) is Sino-Japanese (kango): two Chinese characters read with their on'yomi, readings borrowed from Chinese. They are \emph{tei}, \textquotedblleft{}garden\textquotedblright{} (Middle Chinese \emph{deng}) and \emph{en}, \textquotedblleft{}park\textquotedblright{} (\emph{hjwon}) (KANJIDIC2; Middle Chinese in Baxter's transcription).
\end{cousinweb}

\begin{grammarlens}[title={What you've built}]
One more word for getting around.
\end{grammarlens}

\subsection*{Your turn}
\begin{itemize}
\raggedright
  \item \emph{Say it:} \emph{teien}
  \item \emph{Say it:} \emph{teien}, once more
  \item \emph{Say it:} \emph{tera no teien wa shizuka deshita} — The temple garden was quiet.
  \item \emph{Recall:} say the Japanese for wonderful, then the Japanese for to look down on, then say \emph{teien} again
\end{itemize}

\begin{tcolorbox}[breakable,colback=teal!4,colframe=teal!35!black,title={Before you move on}]
What does \ja{ていえん} mean? (\textquotedblleft{}A landscaped garden\textquotedblright{}.) Say it once more.
\end{tcolorbox}
\section[\texorpdfstring{\ja{ゆうめい} (yūmei)}{ゆうめい (yūmei)}]{\ja{ゆうめい} (yūmei) — famous}

\label{lesson:JA-C405-yuumei}



\begin{quote}
Before the new one: say the Japanese for to look down on, then the Japanese for a landscaped garden.
\end{quote}

\subsection*{You'll want to know: \ja{ゆうめい}}
\textbf{\ja{ゆうめい}} — \emph{yūmei} — \textquotedblleft{}famous\textquotedblright{}.

\begin{quote}
\textbf{\ja{この} \ja{ていえんは} \ja{ゆうめいです。}} — \emph{kono teien wa yūmei desu} — \textquotedblleft{}This garden is famous.\textquotedblright{}
\end{quote}

A na-adjective: \ja{な} before a noun.

\begin{cousinweb}
\textbf{\ja{ゆうめい}} (\emph{yūmei}) is Sino-Japanese (kango): two Chinese characters read with their on'yomi, readings borrowed from Chinese. They are \emph{yū}, \textquotedblleft{}have\textquotedblright{} (Middle Chinese \emph{hjuwX}) and \emph{mei}, \textquotedblleft{}name\textquotedblright{} (\emph{mjieng}) (KANJIDIC2; Middle Chinese in Baxter's transcription).
\end{cousinweb}

\begin{grammarlens}[title={What you've built}]
One more word for getting around.
\end{grammarlens}

\subsection*{Your turn}
\begin{itemize}
\raggedright
  \item \emph{Say it:} \emph{yūmei}
  \item \emph{Say it:} \emph{yūmei}, once more
  \item \emph{Say it:} \emph{kono teien wa yūmei desu} — This garden is famous.
  \item \emph{Recall:} say the Japanese for to look down on, then the Japanese for a landscaped garden, then say \emph{yūmei} again
\end{itemize}

\begin{tcolorbox}[breakable,colback=teal!4,colframe=teal!35!black,title={Before you move on}]
What does \ja{ゆうめい} mean? (\textquotedblleft{}Famous\textquotedblright{}.) Say it once more.
\end{tcolorbox}
\section[\texorpdfstring{\ja{みごと} (migoto)}{みごと (migoto)}]{\ja{みごと} (migoto) — splendid; magnificent}

\label{lesson:JA-C405-migoto}



\begin{quote}
Before the new one: say the Japanese for a landscaped garden, then the Japanese for famous.
\end{quote}

\subsection*{You'll want to know: \ja{みごと}}
\textbf{\ja{みごと}} — \emph{migoto} — \textquotedblleft{}splendid; magnificent\textquotedblright{}.

\begin{quote}
\textbf{\ja{さくらが} \ja{みごとに} \ja{さいています。}} — \emph{sakura ga migoto ni saite imasu} — \textquotedblleft{}The cherry trees are in splendid bloom.\textquotedblright{}
\end{quote}

A na-adjective: \ja{な} before a noun.

\begin{cousinweb}
\textbf{\ja{みごと}} (\emph{migoto}) is a native Japanese word (wago), inherited rather than borrowed from Chinese, built from \emph{mi}, \textquotedblleft{}seeing\textquotedblright{} (from \emph{miru}) and \emph{goto}, \textquotedblleft{}thing\textquotedblright{}; its kanji were attached later as kun'yomi (JMdict).
\end{cousinweb}

\begin{grammarlens}[title={What you've built}]
One more word for getting around.
\end{grammarlens}

\subsection*{Your turn}
\begin{itemize}
\raggedright
  \item \emph{Say it:} \emph{migoto}
  \item \emph{Say it:} \emph{migoto}, once more
  \item \emph{Say it:} \emph{sakura ga migoto ni saite imasu} — The cherry trees are in splendid bloom.
  \item \emph{Recall:} say the Japanese for a landscaped garden, then the Japanese for famous, then say \emph{migoto} again
\end{itemize}

\begin{tcolorbox}[breakable,colback=teal!4,colframe=teal!35!black,title={Before you move on}]
What does \ja{みごと} mean? (\textquotedblleft{}Splendid; magnificent\textquotedblright{}.) Say it once more.
\end{tcolorbox}
\section[\texorpdfstring{\ja{あれる} (areru)}{あれる (areru)}]{\ja{あれる} (areru) — to become rough (of the sea or weather); to fall into ruin}

\label{lesson:JA-C405-areru}



\begin{quote}
Before the new one: say the Japanese for famous, then the Japanese for splendid.
\end{quote}

\subsection*{You'll want to know: \ja{あれる}}
\textbf{\ja{あれる}} — \emph{areru} — \textquotedblleft{}to become rough (of the sea or weather); to fall into ruin\textquotedblright{}.

\begin{quote}
\textbf{\ja{たいふうの} \ja{あとで、ていえんが} \ja{あれていました。}} — \emph{taifū no ato de, teien ga arete imashita} — \textquotedblleft{}After the typhoon, the garden was in a mess.\textquotedblright{}
\end{quote}

Polite \textbf{\ja{あれます}}, \emph{aremasu}; an ichidan verb. Past (polite): \textbf{\ja{たいふうで} \ja{うみが} \ja{あれました。}}, \emph{taifū de umi ga aremashita}.

\begin{cousinweb}
\textbf{\ja{あれる}} (\emph{areru}) is a native Japanese word (wago), inherited rather than borrowed from Chinese; its kanji was attached later as a kun'yomi (JMdict).
\end{cousinweb}

\begin{grammarlens}[title={What you've built}]
One more word for getting around.
\end{grammarlens}

\subsection*{Your turn}
\begin{itemize}
\raggedright
  \item \emph{Say it:} \emph{areru}
  \item \emph{Say it:} \emph{areru}, once more
  \item \emph{Say it:} \emph{taifū no ato de, teien ga arete imashita} — After the typhoon, the garden was in a mess.
  \item \emph{Recall:} say the Japanese for famous, then the Japanese for splendid, then say \emph{areru} again
\end{itemize}

\begin{tcolorbox}[breakable,colback=teal!4,colframe=teal!35!black,title={Before you move on}]
What does \ja{あれる} mean? (\textquotedblleft{}To become rough (of the sea or weather); to fall into ruin\textquotedblright{}.) Say it once more.
\end{tcolorbox}
\section[\texorpdfstring{\ja{すずむ} (suzumu)}{すずむ (suzumu)}]{\ja{すずむ} (suzumu) — to cool off (in the shade or a breeze)}

\label{lesson:JA-C405-suzumu}



\begin{quote}
Before the new one: say the Japanese for splendid, then the Japanese for to become rough.
\end{quote}

\subsection*{You'll want to know: \ja{すずむ}}
\textbf{\ja{すずむ}} — \emph{suzumu} — \textquotedblleft{}to cool off (in the shade or a breeze)\textquotedblright{}.

\begin{quote}
\textbf{\ja{ていえんの} \ja{えんがわで} \ja{すずみました。}} — \emph{teien no engawa de suzumimashita} — \textquotedblleft{}I cooled off on the veranda by the garden.\textquotedblright{}
\end{quote}

Polite \textbf{\ja{すずみます}}, \emph{suzumimasu}; a godan verb. Past (polite): \textbf{\ja{おじいさんが} \ja{えんがわで} \ja{すずみました。}}, \emph{ojiisan ga engawa de suzumimashita}.

\begin{cousinweb}
\textbf{\ja{すずむ}} (\emph{suzumu}) is a native Japanese word (wago), inherited rather than borrowed from Chinese; its kanji was attached later as a kun'yomi (JMdict).
\end{cousinweb}

\begin{grammarlens}[title={What you've built}]
One more word for getting around.
\end{grammarlens}

\subsection*{Your turn}
\begin{itemize}
\raggedright
  \item \emph{Say it:} \emph{suzumu}
  \item \emph{Say it:} \emph{suzumu}, once more
  \item \emph{Say it:} \emph{teien no engawa de suzumimashita} — I cooled off on the veranda by the garden.
  \item \emph{Recall:} say the Japanese for splendid, then the Japanese for to become rough, then say \emph{suzumu} again
\end{itemize}

\begin{tcolorbox}[breakable,colback=teal!4,colframe=teal!35!black,title={Before you move on}]
What does \ja{すずむ} mean? (\textquotedblleft{}To cool off (in the shade or a breeze)\textquotedblright{}.) Say it once more.
\end{tcolorbox}
